\clearpage
\fancyhead[L]{\small ANZEIGEN.}
\SourceLeaf{100}
\pdfbookmark[0]{Anzeigen}{anzeigen}
\begin{center}
{\Large\bfseries ANZEIGEN.}\par
\vspace{1.8em}
\rule{.88\textwidth}{.3pt}\par
Göttingische gelehrte Anzeigen.\quad 1821 Februar 26.\par
\rule{.88\textwidth}{.3pt}
\end{center}
\pdfbookmark[1]{1821 Februar 26}{anzeige-1821-02-26}

Am 15. Februar wurde der Königl. Societät vom Hrn Hofr. \textsc{Gauss} eine Vorlesung übergeben, überschrieben
\begin{center}
\textit{Theoria Combinationis observationum erroribus minimis obnoxiae, pars prior,}
\end{center}
\noindent die eine der wichtigsten Anwendungen der Wahrscheinlichkeitsrechnung zum Gegenstande hat. Alle Beobachtungen, die sich auf Grössenbestimmungen aus der Sinnenwelt beziehen, können, mit welcher Genauigkeit und mit wie vortrefflichen Werkzeugen sie auch angestellt werden, nie \textit{absolute} Genauigkeit haben; sie bleiben immer nur Näherungen, grössern oder kleinern Fehlern ausgesetzt. Nicht von solchen Fehlern ist hier die Rede, deren Quellen genau bekannt sind, und deren Grösse bei bestimmten Beobachtungen jedesmal berechnet werden kann; denn da dergleichen Fehler bei den beobachteten Grössen in Abzug gebracht werden können und sollen, so ist es dasselbe, als ob sie gar nicht da wären. Ganz anders verhält es sich dagegen mit den als zufällig zu betrachtenden Fehlern, die aus der beschränkten Schärfe der Sinne, aus mancherlei unvermeidlichen und keiner Regel folgenden Unvollkommenheiten der Instrumente, und aus mancherlei regellos (wenigstens für uns) wirkenden Störungen durch äussere Umstände (z. B. das Wallen der Atmosphäre beim Sehen, Mangel absoluter Festigkeit beim Aufstellen der Instrumente) herrühren. Diese zufälligen Fehler, die
\SourceLeaf{101}
dem Calcül nicht unterworfen werden können, lassen sich nicht \textit{wegschaffen}, und der Beobachter kann sie durch sorgfältige Aufmerksamkeit und durch Vervielfältigung der Beobachtungen nur \textit{vermindern}: allein nachdem der Beobachter das seinige gethan hat, ist es an dem Geometer, die Unsicherheit der Beobachtungen und der durch Rechnung daraus abgeleiteten Grössen nach streng mathematischen Principien zu würdigen, und was das wichtigste ist, da, wo die mit den Beobachtungen zusammenhängenden Grössen aus denselben durch verschiedene Combinationen abgeleitet werden können, diejenige Art vorzuschreiben, wobei so wenig Unsicherheit als möglich zu befürchten bleibt.

Obgleich die zufälligen Fehler als solche keinem Gesetze folgen, sondern ohne Ordnung in einer Beobachtung grösser, in einer andern kleiner ausfallen, so ist doch gewiss, dass bei einer bestimmten Beobachtungsart, auch die Individualität des Beobachters und seiner Werkzeuge als bestimmt betrachtet, die aus jeder einfachen Fehlerquelle fliessenden Fehler nicht bloss in gewissen Grenzen eingeschlossen sind, sondern dass auch alle möglichen Fehler zwischen diesen Grenzen ihre bestimmte relative Wahrscheinlichkeit haben, der zu Folge sie nach Maassgabe ihrer Grösse häufiger oder seltener zu erwarten sind, und derjenige, der eine genaue und vollständige Einsicht in die Beschaffenheit einer solchen Fehlerquelle hätte, würde diese Grenzen und den Zusammenhang zwischen der Wahrscheinlichkeit der einzelnen Fehler und ihrer Grösse, zu bestimmen im Stande sein, auf eine ähnliche Weise, wie sich bei Glücksspielen, so bald man ihre Regeln kennt, die Grenzen der möglichen Gewinne und Verluste, und deren relative Wahrscheinlichkeiten berechnen lassen. Dasselbe gilt auch von dem aus dem Zusammenwirken der einfachen Fehlerquellen entspringenden Totalfehler. Auch sind diese Begriffe nicht auf unmittelbare Beobachtungen beschränkt, sondern auch auf mittelbare aus Beobachtungen abgeleitete Grössenbestimmungen anwendbar. In der Wirklichkeit werden uns freilich fast allemal die Mittel fehlen, das Gesetz der Wahrscheinlichkeiten der Fehler \textit{a priori} anzugeben.

Wie wir die Unzulässigkeit einer bestimmten Art von Beobachtungen im Allgemeinen abschätzen wollen, hängt zum Theil von unserer Willkür ab. Man kann dabei. entweder bloss die Grösse der äussersten möglichen Fehler zum Maassstabe wählen, oder zugleich auf die grössere oder geringere Wahrscheinlichkeit der einzelnen möglichen Fehler mit Rücksicht nehmen. Das letztere scheint angemessener zu sein. Allein diese Berücksichtigung kann auf vielfache Weise ge%
